\documentclass[12pt,a4paper]{article}
\usepackage[T2A]{fontenc}
\usepackage{babel}
\babelprovide[import,main]{russian}
\usepackage[margin=2cm]{geometry}
\linespread{0.91}

\begin{document}

\begin{center}
  {\Large\bfseries Fixed Set: совершенное хеширование FKS}
\end{center}

\section*{Устройство и построение}
\begin{itemize}
  \setlength{\topsep}{3pt}
  \setlength{\itemsep}{0pt}
  \item \textbf{Первый уровень.} Распределяем $n>0$ различных ключей по $n$ временным корзинам.
    Если сумма квадратов их размеров $\sum_i s_i^2>4n$, выбираем новый хеш.
  \item \textbf{Второй уровень.} Для корзины размера $s_i$ выделяем $s_i^2$ ячеек \texttt{std::optional<int>}.
    При коллизии очищаем ячейки и выбираем новый хеш.
    Храним готовые таблицы и хеши обоих уровней, но не временные корзины.
  \item \textbf{Хеширование.} Индекс равен $((ax+b)\bmod p)\bmod M$, где $M$ --- размер таблицы,
    а $p=2\,000\,000\,011$ --- простое.
    Выбираем $a$ из $[1,p-1]$ и $b$ из $[0,p-1]$ независимо и равномерно.
  \item К отрицательному остатку ключа прибавляем $p$.
    Для разных ключей из $[-10^9,10^9]$ выполняется $0<|x-y|\le2\cdot10^9<p$, поэтому их вычеты различны.
    Для пустой таблицы индекс не вычисляем.
\end{itemize}

\section*{Корректность}
\begin{itemize}
  \setlength{\topsep}{3pt}
  \setlength{\itemsep}{0pt}
  \item Каждый ключ попадает в одну корзину, а вторичная таблица принимается только после записи всех её ключей
    без коллизий.
  \item Поиск теми же хешами находит ячейку ключа.
    Пустая ячейка или несовпадение ключей дают \texttt{false}.
  \item Повторный \texttt{initialize} уничтожает старые таблицы и строит новые только из нового входа.
    Пустой вход оставляет множество пустым.
\end{itemize}

\section*{Время}
\begin{itemize}
  \setlength{\topsep}{3pt}
  \setlength{\itemsep}{0pt}
  \item \textbf{Модель.} Вход фиксирован, пары $(a,b)$ выбираются независимо и равномерно в каждой попытке.
    В коде выбор имитирует \texttt{std::mt19937\_64} с фиксированным начальным состоянием.
    По универсальности семейства вероятность коллизии пары не выше $1/M$.
  \item \textbf{Первый уровень.} Для числа столкнувшихся пар $C$ имеем $\sum_i s_i^2=n+2C$ и
    $\mathrm{E}C\le(n-1)/2$.
    По неравенству Маркова вероятность превысить $4n$ меньше $1/2$.
    В среднем нужно не более двух попыток по $O(n)$ с учётом очистки.
  \item \textbf{Второй уровень.} Для непустой корзины вероятность любой коллизии не выше
    $s_i(s_i-1)/(2s_i^2)<1/2$ по сумме вероятностей для пар.
    В среднем нужно не более двух попыток по $O(s_i^2)$ с учётом очистки.
    Пустая корзина требует $O(1)$ времени.
  \item По условию $\sum_i s_i^2\le4n$ построение занимает ожидаемое $O(n)$ при $n>0$.
    Запрос занимает $O(1)$ в худшем случае, а повторная инициализация добавляет до $O(n_{\mathrm{old}})$ на очистку.
\end{itemize}

\section*{Память}
\begin{itemize}
  \setlength{\topsep}{3pt}
  \setlength{\itemsep}{0pt}
  \item Временные корзины и готовые таблицы занимают $O(n+1)$ с учётом ёмкости векторов.
    Вторичные ячейки выделяются только после проверки $\sum_i s_i^2\le4n$.
    Буферы неудачных попыток освобождаются или используются снова.
  \item При повторной инициализации пик --- $O(n_{\mathrm{old}}+n+1)$:
    старые таблицы освобождаются после распределения ключей, но до выделения новых вторичных таблиц.
    После возврата остаётся $O(n+1)$ без учёта входного вектора.
\end{itemize}

\end{document}
